\begin{abstract}
  We develop a local thermodynamic interpretation of the projective spectral entropy
  functional $\Spi[g] = \frac{1}{2}\log\det' A_g$, extending the covariant framework
  established in a companion paper.
  The diagonal heat kernel $K(x,x;t)$ defines a local spectral energy density
  $u(x;t) = -\partial_t \log K(x,x;t)$, whose Seeley--DeWitt expansion yields
  a local spectral multiplier field $\beta(x)^{-1} = \beta_0^{-1} - \frac{1}{6}R(x) + O(\nabla^2 R)$.
  Here $\beta_0^{-1} = 2/t_*$ is a universal UV contribution fixed by the spectral
  admissibility scale $t_*$, while the curvature correction is geometric and covariant.

  We then formulate a local spectral first law as a compatibility theorem:
  in the infrared-dominant regime, the metric variation of $\Spi$ and the metric
  variation of the projected matter functional are related through the local multiplier
  field, yielding a constraint $\delta\Spi = \int \beta(x)^{-1}\,\delta E_\Pi(x)$.
  This is not a definition but a structural consequence of the Seeley--DeWitt
  hierarchy and the induced-gravity normalization established in the companion paper.

  We then impose a spectral equilibrium principle: admissible geometries extremize
  $\Spi$ at fixed projected matter content.
  The resulting Euler--Lagrange equation is Einstein's field equation
  $G_{\mu\nu} = 8\pi G\, T_{\mu\nu}$, with the factor $8\pi G$ arising purely from
  the normalization conventions of the geometric and matter sides.
  The derivation requires no horizon structure, no Rindler wedge, no Raychaudhuri
  theorem, and no assumption of local Lorentz invariance at the Planck scale.
  Einstein's equation emerges as a spectral extremum condition within the
  effective infrared geometry.
\end{abstract}
